% Session additions to the implementation reference only.
\clearpage
\subsection{Tab controllers, reusable logging and retained index searches}
\label{sec:index-gui}
\code{gui/n2_db.py} owns the main window, settings and coordination between
tabs. \code{AnomalyTabController} in \code{gui/anomaly_tab.py} owns anomaly
input loading, Build/Stop, process output and build state. The parallel
\code{IndexTabController} in \code{gui/index_tab.py} owns index search,
selection, calculation and retained results. Both are \code{QObject}s;
their action methods are connected to widgets loaded from \code{gui/n2_db.ui}.

The index tab has two side-by-side areas. The left contains
\textbf{Search theories with empty indices}, \textbf{Select all}, and
checkboxes grouped by the theory's ordered gauge-group factors, with counts.
The right contains a read-only log and \textbf{Calculate index}, which becomes
\textbf{Stop} during calculation. Select all checks or clears every listed
group; individual selections update its state too.

\code{gui/index_search.py} runs as a separate Sage process. It uses the saved
MySQL settings and \code{iter_lagrangian_index_jobs}, without schema migration
or writes. An index job is selected when its full index, Coulomb index
\emph{or} Coulomb spectrum is SQL NULL or JSON null. A complete record of
lower or unknown precision is not an ``empty'' result and is not selected by
this GUI search. The separate CLI's upgrade mode remains available.
One Lagrangian realization is retrieved for each shared theory.

Successful searches retain complete theory/realization IDs, factor and hyper
inputs, missing-field metadata and source identity in the controller.
\code{retrieved_jobs} returns all jobs and \code{selected_jobs} returns checked
groups' jobs as copies. The next calculation uses this snapshot without a
second database-wide search. New groups begin unchecked; a successful refresh
preserves checks for still-present groups. An empty successful search clears
the list; a failed/incomplete search retains the last successful snapshot.
Changing the source database clears it; changing a cutoff does not. The
snapshot is in memory only and does not survive restarting the application.

\textbf{Shared log implementation.} \code{gui/logging_utils.py} provides
\code{redact_message}, \code{make_log_record}, \code{format_log_message},
and \code{GuiLogger}. A logger accepts an append callback, secrets to redact,
and an optional file started with a configurable prefix. Both tab controllers
use this API, so future logs can reuse the same timestamp, level, multiline,
redaction and flushing behavior. Each index search and calculation attempt
creates its own file, including failures before worker startup:
\begin{lstlisting}
logs/log_index_YYYYMMDD_HHMMSS_ffffff.log
\end{lstlisting}
The file and read-only widget receive the same redacted messages. A file error
is reported on screen without stopping the work. The anomaly prefix remains
\code{log_anomalies}. Files close after completion, failure or cancellation.

\clearpage
\subsection{Multiprocessing index calculation and partial progress}
\code{gui/index_calculator.py} is the Sage coordinator started by
\textbf{Calculate index}. It distributes the selected retained theories across
spawned processes. With $C$ selected theories and CPU limit $P$, the pool uses
at most $\min(C,P)$ workers and at most twice that many outstanding tasks.
One task contains one theory; workers receive further tasks dynamically.
The setting $-1$ resolves to all logical CPUs, with a one-worker fallback.

Each worker owns its MySQL connection, opened against the initialized schema,
and closes it at normal shutdown. Connections never cross process boundaries.
The worker calls \code{common.n2_theory_db_indices.calculate_index_job} with
saved full/Coulomb cutoffs, executable paths, timeout and both cache filenames.
Inner full-index calculations use one process to avoid nested worker pools.
The worker passes its own theory connection to the disconnected-sector reuse
path described in Section~\ref{sec:disconnected}.

\begin{enumerate}
\item Recheck the retained theory/realization IDs for fields currently missing.
This is a targeted stale-job check, not another search of all empty theories.
\item Compute missing full index, Coulomb index and complete Coulomb spectrum
outside write transactions. A stored full index does not limit the separate
Coulomb calculation's requested precision.
\item Save each successful component independently using
\code{missing_only=True}. A concurrent fill, sufficient result or legacy
unknown-precision value is preserved. A failure in one component does not
erase successful components already committed.
\item Publish outcomes. Confirmed complete jobs are removed from the retained
selection when the run ends. Failed, stopped or unconfirmed jobs remain for
retry; a later retry rechecks which components still need work.
\end{enumerate}

A manager-backed queue streams log records from workers to the coordinator.
The coordinator drains results and logs while tasks are active, avoiding the
feeder-thread shutdown pattern in which a child waits on a full, undrained
ordinary multiprocessing queue. Scheduling remains bounded even for a long
selection. Unexpected worker exit is reported rather than silently retried.

\textbf{Stop and close.} Stop cancels queued work where possible and lets
active components finish and save before workers shut down. Window close
requests the same cooperative stop. Search, anomaly Build and Preferences
are mutually blocked during a calculation. A read-only search can be cancelled
on close. These ownership and shutdown rules address process/queue deadlocks;
they do not eliminate mathematical workloads that take a long time.

\clearpage
\subsection{Database lock ordering and the test-permission repair}
The index writer and legacy-cutoff recorder lock the parent theory before
property/realization rows. Existing-theory anomaly attachment uses the same
parent-first order. This removes the earlier child/parent cycle between
imports and index updates. The fresh-import missing-row gap-lock fix remains:
a newly inserted theory gets a direct property insert rather than a locking
lookup for an absent property row.

For MySQL errors 1205 (lock wait timeout) and 1213 (deadlock), the index write
and cutoff APIs roll back the whole short transaction, wait 50 then 100 ms,
and retry at most twice. Each attempt reads the current state again and merges
only permitted fields; the expensive index is not recomputed. Connection
failures and failed rollbacks are not replayed. Other clients can still
introduce contention, so the result is bounded recovery, not a claim that all
possible database deadlocks are impossible.

The 14 September recorded index run passed 338 tests in 55.104 s, including
the real GUI search/calculation protocol, custom caches, partial failures,
cooperative Stop, abrupt worker exit, connection cleanup, an eight-worker
256-theory stress case and targeted parent-row contention. The controlled
stress and contention runs added no observed deadlocks.

The later \code{all_test.sh} failure was a test-account permission problem:
concurrency tests queried server-wide \code{INNODB_METRICS}, which requires
privileges the database-scoped test account does not have. The repaired tests
observe rollback calls on their own connections, including successfully retried
lock failures. Cleanup inspects only connections owned by that test user.
No global PROCESS or SUPER grant is required. The recorded post-repair suite
passed 339 tests. Direct-file/module cache-builder resume coverage remains in
\code{test/test_build_decomposition_cache.py}; its appearance at the start of
the submitted test log did not identify it as the failing test.

\textbf{Test organization.} All new tests and test helpers live in \code{test/}.
The moved GUI tests use \code{test_gui_} names; the opt-in real credential-vault
check is \code{test/check_secret_service.py}. Application files under
\code{gui/} contain runtime code only. \code{all_test.sh} discovers the combined
suite from \code{test/}. Live integration fixtures reset only the configured
test database and need database-scoped privileges; ordinary GUI tests use
temporary settings and fixture credentials.

\clearpage
\subsection{Password-confirmed deletion from Settings}
\label{sec:delete-settings}
The MySQL section of Settings has \textbf{Delete all database contents...}.
\code{SettingsDialog} snapshots the displayed database, host, port, user and
timeout, including unsaved edits, plus the effective Unix socket when set.
\code{DeleteDatabaseDialog} shows that target, an irreversible-deletion warning,
a fresh masked password field, and Delete/Cancel. The saved password is not
prefilled. Opening or cancelling the dialog performs no database operation.

The editable layout is \code{gui/database_clear_dialog.ui}. Python loads it
with \code{uic.loadUi}, as it loads the main and Settings forms. The UI file
owns labels, layout, password echo mode, disabled initial Delete button,
default Cancel button and tab order. \code{gui/database_clear_dialog.py}
supplies the target text and platform warning icon, connects signals, and
controls the asynchronous operation. It no longer constructs the layout.

\begin{enumerate}
\item A nonempty freshly entered password enables Delete. Password whitespace
is preserved. Delete starts a separate Sage process for
\code{gui.database_clear}; target, password and explicit confirmation travel
through stdin, never through command-line arguments or saved preferences.
\item The worker authenticates a fresh MySQL connection using the entered
password and \code{initialize_schema=False}. This is database authentication,
not comparison with the previously saved credential.
\item \code{delete_all_theory_data} checks schema version 7, starts a transaction,
executes \texttt{DELETE FROM theories}, and commits. Foreign-key cascades remove
the associated property, realization, gauge, hyper, flavor, index and spectrum
data. Tables, schema metadata, account and SQLite caches remain. There is no
DROP, TRUNCATE or auto-increment reset. Exceptions trigger rollback.
\item A confirmed result returns the deleted parent-theory count and clears
the index tab's retained jobs, even if Settings is subsequently cancelled.
\end{enumerate}

The dialog remains responsive but disables password input and cancellation
during deletion. Dialog, Settings and main-window close wait for the worker.
Settings is unavailable during a build, search or calculation. Wrong-password
rejection allows a fresh entry; errors and logs redact secrets. A lost worker
or connection can make the commit outcome uncertain, so the UI reports an
unconfirmed result and advises inspecting the database before retrying.
The helper guarantees the transactional attempt, not knowledge of a commit
after a lost connection.

\textbf{Recorded validation, 16 September.} The full suite passed 370 tests in
65.584 s with no skips; an additional actual stdin-worker authentication test
passed afterward in 1.444 s. After extracting the layout, 38 deletion-dialog
and main-GUI tests passed in 1.477 s and the dialog was rendered for visual
inspection. Live deletion fixtures checked every theory-data table, wrong
password preservation, rollback after a failure before commit, and reuse of
the preserved schema. These counts describe separate runs, not a claimed
371-test combined run. No production database was cleared.
